\documentclass[11pt]{article}
\usepackage[margin=1in,top=.85in,bottom=.85in]{geometry}
\usepackage{amsmath,amssymb,amsthm,hyperref}
\hypersetup{colorlinks=true,linkcolor=blue,urlcolor=blue,citecolor=blue}
\setlength{\emergencystretch}{2em}
\newtheorem{theorem}{Theorem}
\newtheorem{lemma}[theorem]{Lemma}
\title{Classical-mean power sums:\\the exact real-exponent ranges}
\author{Research proof and formalization packet}
\date{October 26, 2026}
\begin{document}
\maketitle
\begin{abstract}
For the arithmetic, geometric, harmonic and quadratic means of two
positive real numbers, we determine every real exponent for which
$A^n+G^n\leq Q^n+H^n$ holds universally. The range is exactly
$n\leq0$ or $n\geq1$. The reverse inequality holds universally only
at $n=0$. Consequently both strict signs occur for each $0<n<1$.
The first classification contradicts the proposed range in
Me\v{s}trovi\'c--Andji\'c's Conjecture~1.5 in its published version.
Classical two-coordinate convexity proves sufficiency; a negative
second derivative at equal inputs proves necessity. Explicit positive
inputs establish failure of the reverse at every nonzero exponent.
A complete five-module Lean certificate verifies both classifications.
The result is one combined source-question resolution; known integer
cases and classical convexity tools are credited.
\end{abstract}

\section{The source question and exact result}
For positive real numbers $x,y$, write
\[
 H=\frac{2xy}{x+y},\qquad G=\sqrt{xy},\qquad
 A=\frac{x+y}{2},\qquad Q=\sqrt{\frac{x^2+y^2}{2}}.
\]
Me\v{s}trovi\'c and Andji\'c \cite{MA} prove
$A^n+G^n\leq Q^n+H^n$ for integer $n$. Their published
Conjecture~1.5 proposes validity for every negative real exponent and
every $n\geq1/2$, with neither direction universally valid for
$0<n<1/2$. The preprint numbers the same assertion as
Conjecture~1.6. This is the source question represented by
MathDB~338823.

\begin{theorem}\label{thm:classification}
For a real number $n$, the following equivalences hold:
\begin{align}
 [\,A^n+G^n\leq Q^n+H^n\text{ for all }x,y>0\,]
 &\quad\Longleftrightarrow\quad n\leq0\text{ or }n\geq1,
 \label{eq:forward}\\
 [\,Q^n+H^n\leq A^n+G^n\text{ for all }x,y>0\,]
 &\quad\Longleftrightarrow\quad n=0.
 \label{eq:reverse}
\end{align}
For every $0<n<1$, there are positive inputs for which each strict
direction holds. All powers are ordinary real powers of positive bases.
\end{theorem}

The result addresses an entire interval of the proposed source range:
the first direction fails universally for every $1/2\leq n<1$.
The source's integer results and the equality at zero are known cases,
not additional claimed outcomes. Both equivalences and all witnesses
constitute one classification.

\section{Known mean comparisons and convexity}
We first record elementary relations, including the source's sum and
product comparisons:
\begin{equation}\label{eq:meanrelations}
 H\leq G\leq A\leq Q,\qquad A+G\leq Q+H,\qquad QH\leq AG.
\end{equation}
For completeness, put $c=A/G\geq1$. The identities
$AH=G^2$ and $Q^2=2A^2-G^2$ give
\begin{equation}\label{eq:normalize}
 (H,G,A,Q)/G=(c^{-1},1,c,\sqrt{2c^2-1}).
\end{equation}
Let $q=\sqrt{2c^2-1}$. The order relations are immediate. The remaining
two comparisons follow from
\[
 q^2-(c+1-c^{-1})^2=\frac{(c-1)^3(c+1)}{c^2}\geq0,
 \qquad c^4-q^2=(c^2-1)^2\geq0.
\]
All compared quantities are positive, so these squared comparisons
give $q\geq c+1-c^{-1}$ and $q/c\leq c$. This proves
\eqref{eq:meanrelations}. These source comparisons are credited as
known; the certificate also proves them directly from the definitions.

\begin{lemma}\label{lem:convex}
Suppose $d\leq b\leq a\leq c$ and $a+b\leq c+d$ in an interval.
For every increasing convex function $f$ on that interval,
\[
 f(a)+f(b)\leq f(c)+f(d).
\]
\end{lemma}
\begin{proof}
Put $c'=a+b-d$. Then $d\leq b\leq a\leq c'\leq c$ and
$a+b=c'+d$. If $c'=d$, the equal-sum comparison is equality.
Otherwise, applying the chord inequality on $[d,c']$ to $a$ and $b$
and adding gives $f(a)+f(b)\leq f(c')+f(d)$. Monotonicity gives
$f(c')\leq f(c)$, proving the assertion.
\end{proof}
This is the classical two-coordinate weak-majorization argument,
included to make the proof self-contained; it is not a new convexity
principle.

For $n\geq1$, apply Lemma~\ref{lem:convex} to $f(t)=t^n$ and
$(a,b,c,d)=(A,G,Q,H)$, using \eqref{eq:meanrelations}. This proves
the first inequality throughout that range.

For $n\leq0$, the four numbers
\[
 d=n\log Q\leq b=n\log A\leq a=n\log G\leq c=n\log H
\]
have $a+b\leq c+d$. Indeed $AG\geq QH$, and multiplication of
their logarithmic comparison by $n\leq0$ reverses its direction.
Apply the lemma to the increasing convex function $\exp$ and reorder
the terms. This proves the first inequality for every $n\leq0$,
including equality at zero.

\section{Failure for every intermediate exponent}
Every $c\geq1$ in \eqref{eq:normalize} is realized by the actual
positive pair
\begin{equation}\label{eq:pair}
 x_c=c+\sqrt{c^2-1},\qquad y_c=c-\sqrt{c^2-1}.
\end{equation}
Their product is one, their sum is $2c$, and
$\sqrt{c^2-1}<c$, so both inputs are strictly positive.
Define the normalized defect
\begin{equation}\label{eq:defect}
 D_n(c)=(2c^2-1)^{n/2}+c^{-n}-c^n-1.
\end{equation}
It is precisely $Q^n+H^n-A^n-G^n$ on the pair \eqref{eq:pair}.
Real powers are smooth in a neighborhood of $c=1$, since both bases
are positive there. Direct differentiation gives
\begin{align*}
 D_n(1)&=0,\\
 D'_n(c)&=2nc(2c^2-1)^{n/2-1}
              -nc^{-n-1}-nc^{n-1},\\
 D'_n(1)&=0,\qquad D''_n(1)=4n(n-1).
\end{align*}
For every $0<n<1$, the last value is negative. Since
$D'_n(1)=0$, differentiability of $D'_n$ at one implies
$D'_n(c)<0$ for all $c>1$ sufficiently close to one.
The mean value theorem therefore gives $D_n(c)<D_n(1)=0$
on such inputs. Thus every exponent in $(0,1)$ fails the forward
universal inequality. Together with sufficiency this proves
\eqref{eq:forward}.

An exact counterexample to the source's proposed half-power case is
\[
 n=\tfrac12,\qquad x=4+\sqrt{15},\qquad y=4-\sqrt{15}.
\]
Here $(A,G,H,Q)=(4,1,1/4,\sqrt{31})$, and the proposed inequality
would read $3\leq31^{1/4}+1/2$. In fact $31<(5/2)^4$, so its
right-hand side is strictly smaller than three. This witness is one
instance of the same classification, not another outcome.

\section{The reverse universal range}
We provide strict witnesses against the reverse for every $n\ne0$.
First suppose $n>0$. Put
\[
 \varepsilon=(5/4)^n-1>0,\qquad
 c=\max\{2,(2/\varepsilon)^{1/n}\}.
\]
Then $\varepsilon c^n\geq2$. For $c\geq2$,
$q=\sqrt{2c^2-1}\geq(5/4)c$, since
$(2c^2-1)-(25/16)c^2=(7/16)c^2-1>0$.
Consequently
\[
 Q^n=q^n\geq(5/4)^nc^n
       =(1+\varepsilon)c^n\geq c^n+2>A^n+G^n
\]
on the positive pair \eqref{eq:pair}. Since $H^n>0$, the forward
inequality is strict there and the reverse fails.

If $n<0$, instead put $c=\max\{2,3^{1/(-n)}\}$. Then
$H^n=c^{-n}\geq3$, while $A^n=c^n\leq1$ and $G^n=1$.
Again $Q^n+H^n>A^n+G^n$. At $n=0$, both sums are two.
This proves \eqref{eq:reverse}. Combining these strict witnesses
with the preceding negative-defect witnesses proves the assertion
about both signs throughout $(0,1)$.

\clearpage
\section{Formal coverage and prior-art limits}
The five modules \texttt{MeanAlgebra}, \texttt{MeanSufficiency},
\texttt{MeanNecessity}, \texttt{MeanClassification}, and
\texttt{MeanConverse} contain 34 local declarations. They begin with
the literal positive-input mean definitions. The certificate proves
all algebraic comparisons and normalized input identities, the
two-coordinate convexity lemma, sufficiency, the derivative and
neighborhood-sign argument, actual positive-input counterexamples,
the exact half-power example, and the reverse-direction witnesses.
The endpoint \path{source_full_classification}, in namespace
\texttt{ClassicalMeanRange}, states both equivalences without any
supplied mean, derivative, normalization or witness premise.

A fresh isolated replay compiled every declaration with warnings as
errors under Lean~4.34.1 and the pinned mathlib dependencies. Compiler
and mathlib commits, respectively, are
\begin{center}\small\ttfamily
5045d0056413266e57c625dcd7c365b10e377c52\\
d13f23b723b8a846827a245b89c10fc7d3f11612.
\end{center}
All declaration types and axiom reports are retained. Only
\texttt{propext}, \texttt{Classical.choice}, and \texttt{Quot.sound}
occur; there are no proof holes, custom axioms, unsafe declarations or
native computation assumptions. Exact dependency pins and the replay
script accompany the candidate.

The actual published source retains its conjecture on page~1894.
Its integer theorem and sum/product comparisons, and the classical
convexity principle, receive prior credit. No separate outcome is
claimed for those ingredients, zero-exponent equality, an individual
witness or the reverse direction.

An independent source, mathematical and formal review, including a
fresh replay of all 34 declarations, found no decisive objection.
Targeted literature and posted-solution searches located no earlier
full resolution. Related primary comparisons were also examined:
Me\v{s}trovi\'c's later chain paper \cite{Chain} concerns fixed nested
means; Jameson's Theorems~1--2 \cite{Jameson} compare a power mean
$M_p$ with $(1-p)G+pA$ and with $G^{1-p}A^p$, rather than this
four-mean power-sum family. The complete publisher HTML of the latter
was read even though direct PDF retrieval returned HTTP~403.

The literature assessment is bounded. Subscription citation services
were unavailable, and a broad survey was only partially inspected.
Citation counts and an absence of search hits do not prove priority.
No claim of exhaustive worldwide novelty is made; exact sources,
retrieval limits and independent review are retained in the packet.
No external posting or author contact has occurred.

\begin{thebibliography}{9}
\bibitem{MA}
R.~Me\v{s}trovi\'c and M.~Andji\'c,
\emph{Two Curious Inequalities Involving Different Means of Two
Arguments}, Resonance \textbf{28}(12) (2023), 1891--1897,
\href{https://doi.org/10.1007/s12045-023-1726-x}{doi:10.1007/s12045-023-1726-x};
online January~31, 2024.
Preprint \href{https://arxiv.org/abs/1804.00542v1}{arXiv:1804.00542v1},
Conjecture~1.6; published Conjecture~1.5.
\bibitem{Chain}
R.~Me\v{s}trovi\'c, \emph{A chain of eight inequalities involving means
of two arguments}, The Teaching of Mathematics \textbf{27}(1)
(2024), 27--32,
\href{https://doi.org/10.57016/TM-XDVI2817}{doi:10.57016/TM-XDVI2817}.
\bibitem{Jameson}
G.J.O.~Jameson, \emph{Some inequalities for power means; a problem
from ``The logarithmic mean revisited''}, American Mathematical Monthly
\textbf{130}(3) (2023), 276--278,
\href{https://doi.org/10.1080/00029890.2022.2153560}{doi:10.1080/00029890.2022.2153560}.
\end{thebibliography}
\end{document}
